\chapter{Deep Dive: Configuration --- Config Server, Secrets, and Refresh at Scale}
\label{ch:dd-config}

Chapter~2 introduced the config-server as ``one more property source in
the stack'': a small HTTP service that hands each application its YAML at
startup. That description is complete for a laptop and dangerously
incomplete for anything else. This chapter takes the project's
configuration path apart --- what the server actually serves, who can
read it, what happens when it changes, and what happens when it is not
there --- and then rebuilds it two different ways.

Three questions organize the chapter:

\begin{enumerate}
  \item \textbf{What problem does centralized configuration solve}, and
  which of its three backends (native, git, platform) solves it best for
  which system?
  \item \textbf{What is the blast radius of a configuration file} ---
  who reads it, what precedence rules apply, and what a default value in
  a git-tracked file really means?
  \item \textbf{What breaks} when the server is down, when a value must
  change at runtime, and when a secret must rotate?
\end{enumerate}

\section{Why centralize, and what it solves}

The twelve-factor rule is \emph{strict separation of config from code}:
an image is built once and promoted unchanged through environments; only
the configuration around it differs. Every configuration mechanism is a
way of satisfying that rule, and the config-server is the first one the
Spring ecosystem produced. It solves three concrete problems at once:

\begin{itemize}
  \item \textbf{Duplication.} Five services share Kafka addresses,
  tracing endpoints, logging patterns. One document per service in one
  place beats five copies that drift.
  \item \textbf{Environment awareness.} The same service asks for
  \code{/course-service/prod} or \code{/course-service/default} and gets
  the right document, from the same image.
  \item \textbf{Change without rebuild.} With a git backend, a config
  change is a commit --- reviewed, audited, revertible --- and takes
  effect on the next restart or refresh, with no image rebuilt.
\end{itemize}

The project gets the first benefit fully, the second not at all (there
is one profile, \code{default}), and the third only in principle: with
the \code{native} backend the documents are inside the config-server's
own jar, so changing a value \emph{does} mean a rebuild.

\begin{decision}[three backends, three systems]
\textbf{Native classpath} (this project): config and code in one repo,
one pipeline, one version. Right when one team owns everything and every
config change ships with a code change anyway. \textbf{Git backend}:
config in its own repository with its own history; right when
operations and development are different people, when the same service
runs in many environments, or when an auditor asks ``who changed the
timeout and when''. \textbf{Platform-native} (ConfigMap and Secret, no
server): right when every deployment target is Kubernetes, because the
platform already offers versioned, mountable, RBAC-protected
configuration and one fewer JVM to run. The project sits on the first
and is heading for the third (Section~\ref{sec:config-enhance}); the
second is what most enterprises you interview at actually run.
\end{decision}

\section{Reading the real path}

Follow one property, \code{spring.datasource.url} for course-service,
from the file to the running pod. It starts in the config-server's jar:

\begin{lstlisting}[language=yaml]
# config-server/src/main/resources/configurations/course-service.yml
server:
  port: 8082
spring:
  datasource:
    url: jdbc:postgresql://localhost:5432/ts_course   (*@\co{1}@*)
    username: postgres
    password: root                                     (*@\co{2}@*)
app:
  s3:
    endpoint: ${APP_S3_ENDPOINT:http://localhost:9000}  (*@\co{3}@*)
    secret-key: ${APP_S3_SECRET_KEY:minioadmin}
\end{lstlisting}

The service's packaged \code{application.yml} says where to fetch that
document from:

\begin{lstlisting}[language=yaml]
# course-service/src/main/resources/application.yml
spring:
  application:
    name: course-service                                (*@\co{4}@*)
  config:
    import: optional:configserver:http://localhost:8888 (*@\co{5}@*)
\end{lstlisting}

And the container overrides both the address and the value:

\begin{lstlisting}[language=yaml]
# deploy/k8s/base/course-service.yaml (env excerpt)
- { name: SPRING_CONFIG_IMPORT,
    value: "configserver:http://config-server:8888" }    (*@\co{6}@*)
- { name: SPRING_DATASOURCE_URL,
    value: "jdbc:postgresql://postgres-course:5432/ts_course" }
envFrom:
  - secretRef: { name: postgres-course-app }            (*@\co{7}@*)
\end{lstlisting}

\begin{description}
  \item[\co{1}] A literal, not a placeholder: the served document says
  \code{localhost}. It is correct only on a laptop.
  \item[\co{2}] A database password in cleartext, in a git-tracked
  file, served over HTTP. Dev-only in intent; nothing in the file says
  so.
  \item[\co{3}] The \code{\${VAR:default}} idiom: the config-server
  serves the placeholder \emph{unresolved}; each client resolves it
  against its own environment at startup. This is why the same document
  works in Compose and in Kubernetes --- and why a default in it is a
  value every environment will silently use unless told otherwise.
  \item[\co{4}] The lookup key. The server maps
  \code{course-service} to \code{course-service.yml}.
  \item[\co{5}] \code{optional:} --- if nothing answers on 8888, boot
  anyway from the packaged file. Right for a developer running one
  service in an IDE; wrong for a container.
  \item[\co{6}] The container replaces the whole import line: new host,
  and \emph{no} \code{optional:}. An unreachable config-server is now
  fatal at startup. Compose masks this with
  \code{depends\_on: condition: service\_healthy}; Kubernetes does not,
  and the pod crash-loops until the server is Ready (Chapter~15).
  \item[\co{7}] The password from \co{2} is overridden by the Secret's
  \code{SPRING\_DATASOURCE\_PASSWORD}. The cleartext default is
  therefore harmless in the cluster --- \emph{as long as} the Secret
  exists. Remove it and the service does not fail; it connects with
  \code{root}.
\end{description}

\subsection{Precedence, precisely}

Spring Boot builds one \code{Environment} from an ordered list of
property sources; the first source that defines a key wins. For this
project the order that matters, highest first:

\begin{enumerate}
  \item Command-line arguments and \code{SPRING\_APPLICATION\_JSON}
  (unused here).
  \item OS environment variables --- the Deployment's \code{env} and
  \code{envFrom}, with relaxed binding turning
  \code{SPRING\_DATASOURCE\_URL} into \code{spring.datasource.url}.
  \item The document imported from the config-server.
  \item The packaged \code{application.yml} that did the importing.
\end{enumerate}

Two consequences are easy to get wrong in an interview. First, an
imported document outranks the file that imports it: the config-server
can override anything packaged in the jar, but never anything set by
the platform. Second, a \code{\${VAR:default}} placeholder is
\emph{not} a precedence rule; it is resolved after precedence has
already chosen the document. The env var \code{APP\_S3\_SECRET\_KEY}
wins over the default not because env vars outrank the config-server,
but because the config-server's own text asked for it.

\begin{handson}[read what the server serves]
The config-server has no authentication and Compose publishes its port.
Ask it for course-service's document exactly as the client does:
\begin{lstlisting}[language=bash]
$ curl -s http://localhost:8888/course-service/default | \
    python -m json.tool | head -20
{
    "name": "course-service",
    "profiles": ["default"],
    "label": null,
    "version": null,
    "propertySources": [
        {
            "name": "classpath:/configurations/course-service.yml",
            "source": {
                "server.port": 8082,
                "spring.datasource.url": "jdbc:postgresql://localhost...",
                "spring.datasource.username": "postgres",
                "spring.datasource.password": "root",
\end{lstlisting}
Note \code{version: null} --- the native backend has no commit id to
report. A git backend fills it with the SHA, which is how you answer
``which config is this pod running?'' in production.
\end{handson}

\section{Pros and cons, honestly}

\begin{center}
\begin{tabular}{@{}p{0.3\linewidth}p{0.62\linewidth}@{}}
\toprule
\textbf{Strength} & \textbf{In this project} \\
\midrule
One document per service & Real: five files, zero duplication of Kafka or Zipkin addresses. \\
Environment placeholders & Real: the same files serve Compose and Kubernetes untouched. \\
Startup-time fetch, no runtime dependency & Real: once up, a service never talks to the server again. \\
\midrule
\textbf{Weakness} & \\
\midrule
No authentication & Anyone reaching 8888 reads every document, defaults and passwords included. \\
Native = rebuild to change & A timeout tweak is a config-server image build and rollout. \\
Working defaults for secrets & \code{JWT\_SECRET} and \code{root} have usable fallbacks in git. \\
Ordering constraint & Every pod depends on one pod that has nothing to do after boot. \\
One profile only & \code{default} everywhere; environments differ by env var alone. \\
\bottomrule
\end{tabular}
\end{center}

The last row is worth a sentence. Spring profiles were designed for
exactly the dev/staging/prod split, and the config-server serves
\code{course-service-prod.yml} on top of \code{course-service.yml}
automatically when \code{SPRING\_PROFILES\_ACTIVE=prod}. The project
never uses this; every environmental difference is an env var. That is
a defensible simplification --- one mechanism instead of two --- but it
means the served documents can never say ``this value is for
development only''.

\section{What breaks}

\begin{pitfall}[the signing key is in the repository]
Symptom: none, which is the point. Open the served
\code{api-gateway.yml} and \code{auth-service.yml}:
\begin{lstlisting}[language=yaml]
app:
  jwt:
    secret: ${JWT_SECRET:mySecretKeyThatIsAtLeast256Bits...}
\end{lstlisting}
Cause: the default is a \emph{valid} HS256 key, committed to a public
repository, and \code{create-secrets.sh} defaults \code{JWT\_SECRET} to
the same string. Any deployment that forgets to set the variable signs
tokens with a key the whole internet can read --- and with HS256
(Chapter~28) reading the key means minting tokens for any user. The
fix has two halves: remove the default so a missing variable fails at
startup (\code{\${JWT\_SECRET}} with no colon --- Spring refuses to
resolve it), and treat every \code{\${VAR:value}} in a served document
as a question: ``would I be fine if production ran with this?''
Passwords and keys never pass that test; ports and hostnames usually do.
\end{pitfall}

\begin{pitfall}[refreshed one service, broke all of them]
Symptom: after ``rotating the JWT secret'' with \code{POST
/actuator/refresh} on auth-service, every request through the gateway
returns 401 until the gateway is restarted. Cause: refresh is
\emph{per instance}. Auth-service rebuilt its \code{@RefreshScope}
beans and now signs with the new key; the gateway still verifies with
the old one. Even if both are refreshed, there is no overlap window ---
every token issued before the refresh is invalid the moment the
gateway refreshes. A symmetric key shared by a signer and a verifier
cannot be hot-swapped at all; it needs the \code{kid}-based ceremony
of Chapter~28. The general rule: refresh is for values whose old and new
readings can coexist (a feature flag, a log level, a timeout), not for
values that two parties must agree on.
\end{pitfall}

\begin{pitfall}[healthy but wrong]
Symptom: config-server reports \code{UP}, services start, the wrong
values are in effect. Cause: the config-server's health indicator
checks that the \emph{backend} answers --- for native, that the
classpath directory exists. It says nothing about content. A typo in a
YAML key name (\code{acess-key}) is served happily and ignored happily
by the client, which then uses its default. Neither side logs an error.
The defence is in the client: bind configuration to a
\code{@ConfigurationProperties} record with
\code{@Validated} and \code{@NotBlank} on required fields, so a missing
or misspelled key fails at startup with the property name in the
message.
\end{pitfall}

\section{What breaks at 3\,a.m.}

The config-server pod is evicted at 3\,a.m.; nothing else is.

\begin{itemize}
  \item \textbf{Running services: nothing.} Each fetched its document
  once at boot and holds it in memory. No request path touches the
  config-server. This is the property that justifies a single replica.
  \item \textbf{A service that restarts during the outage: crash loop.}
  Its import is non-optional; the fetch fails; the JVM exits; kubelet
  restarts it with exponential back-off (Chapter~24). When the
  config-server returns, the next restart succeeds --- after up to five
  minutes of back-off delay. The system converges; it does not
  converge quickly.
  \item \textbf{A rolling update during the outage: stalls.} The new
  pod never becomes Ready, the old pod keeps serving (Chapter~15's
  \code{maxUnavailable}), and the rollout waits. Correct behaviour, but
  a deploy pipeline will report a timeout.
  \item \textbf{With \code{optional:} restored: worse.} The restarting
  pod boots from packaged defaults --- \code{localhost:5432},
  \code{localhost:9092} --- becomes Ready (its probe is
  \code{/v3/api-docs}, which needs no database), receives traffic, and
  fails every request. An outage that looked like ``config-server down''
  is now ``course-service returning 500s'', two hops from the cause.
\end{itemize}

Retry with back-off is the middle ground: keep the import mandatory,
but let the client try for a while before dying. It needs
\code{spring-retry} and \code{spring-boot-starter-aop} on the classpath,
then:

\begin{lstlisting}[language=yaml]
spring:
  cloud:
    config:
      fail-fast: true          # unreachable server is fatal ...
      retry:
        max-attempts: 20       # ... after 20 tries
        initial-interval: 2000 # 2 s, then x1.5 each time
        multiplier: 1.5
        max-interval: 15000    # capped at 15 s: ~4 minutes total
\end{lstlisting}

Compare with letting kubelet do the retrying: the JVM restart costs
twenty seconds of class loading each time and the back-off grows to
five minutes; in-process retry costs nothing and caps at fifteen
seconds. Both converge; one wastes a CPU the two-core box does not have.

\section{Enhancing the resolution}
\label{sec:config-enhance}

Two upgrades, in the order an enterprise would do them.

\subsection{Upgrade 1: git backend with encrypted values}

Keep the config-server, move its documents to a repository, and stop
committing secrets in cleartext. The server side:

\begin{lstlisting}[language=yaml]
# config-server/src/main/resources/application.yml
server:
  port: 8888
spring:
  application:
    name: config-server
  profiles:
    active: git                                        (*@\co{1}@*)
  cloud:
    config:
      server:
        git:
          uri: https://git.example.com/ts/config-repo.git
          default-label: main                          (*@\co{2}@*)
          search-paths: '{application}'                (*@\co{3}@*)
          clone-on-start: true                         (*@\co{4}@*)
          username: ${CONFIG_GIT_USER}
          password: ${CONFIG_GIT_TOKEN}
encrypt:
  key: ${CONFIG_ENCRYPT_KEY}                           (*@\co{5}@*)
management:
  endpoints:
    web:
      exposure:
        include: health
\end{lstlisting}

\begin{description}
  \item[\co{1}] The profile selects the backend. \code{native} and
  \code{git} are the two built in; Vault and JDBC exist as well.
  \item[\co{2}] The label is a branch or tag. Pin a production
  environment to \code{release-2026.09} and a config change is a merge,
  not an edit.
  \item[\co{3}] One folder per service in the repository:
  \code{course-service/application.yml}. The same file names as today,
  one directory deeper.
  \item[\co{4}] Clone at startup rather than on first request, so the
  health check proves the credentials work before any client asks.
  \item[\co{5}] A symmetric key for the \code{/encrypt} and
  \code{/decrypt} endpoints. It is the one secret the config-server must
  hold, and it must come from the platform (a Kubernetes Secret), never
  from a file in either repository.
\end{description}

With the key in place, encrypting a value is one request, and the
repository stores only the ciphertext:

\begin{lstlisting}[language=bash]
$ curl -s -X POST http://localhost:8888/encrypt \
       -d 'a-real-256-bit-signing-key-generated-with-openssl'
AQBt3s6...b3f0e   # ciphertext; safe to commit
\end{lstlisting}

\begin{lstlisting}[language=yaml]
# config-repo/api-gateway/application.yml
app:
  jwt:
    secret: '{cipher}AQBt3s6...b3f0e'   # decrypted by the server on serve
\end{lstlisting}

The \code{\{cipher\}} prefix tells the server to decrypt before serving.
The repository now holds no usable secret, the default-value trap is
gone (there is no default), and \code{git log -- api-gateway/} answers
who changed the key and when. What it does \emph{not} fix: the served
document still crosses the network in cleartext to an unauthenticated
endpoint. Add \code{spring-boot-starter-security} to the config-server
with HTTP basic credentials, and give each client
\code{spring.cloud.config.username/password} from its own Secret.

\subsection{Upgrade 2: retire the server}

On Kubernetes the platform can serve the documents itself. Kustomize
generates a ConfigMap per service straight from the existing files, so
the git history is unchanged:

\begin{lstlisting}[language=yaml]
# deploy/k8s/base/kustomization.yaml (excerpt)
configMapGenerator:
  - name: course-service-config
    files:
      # key in the ConfigMap = application.yml, value = the file
      - application.yml=../../config-server/src/main/resources/configurations/course-service.yml
generatorOptions:
  disableNameSuffixHash: false   # hash suffix: a changed file = a
                                 # new ConfigMap name = a rollout
\end{lstlisting}

The Deployment mounts it as a directory and points Spring at the
directory instead of the server:

\begin{lstlisting}[language=yaml]
# deploy/k8s/base/course-service.yaml (excerpt)
spec:
  template:
    spec:
      containers:
        - name: course-service
          env:
            # replaces configserver:http://config-server:8888
            - { name: SPRING_CONFIG_IMPORT, value: "file:/config/" }
            - { name: SPRING_DATASOURCE_URL,
                value: "jdbc:postgresql://postgres-course:5432/ts_course" }
          envFrom:
            - secretRef: { name: postgres-course-app }
          volumeMounts:
            - { name: config, mountPath: /config, readOnly: true }
      volumes:
        - name: config
          configMap: { name: course-service-config }
\end{lstlisting}

Spring Boot treats a trailing slash as a directory import and loads
\code{application.yml} from it, with exactly the precedence the
config-server document had: above the packaged file, below env vars.
Nothing in the service's code or jar changes; Compose keeps working
because the packaged \code{optional:configserver:} import still points
at the (still running, in Compose) server. Then delete
\code{config-server.yaml} from the base, and the cluster has one fewer
JVM, one fewer startup dependency, and one fewer unauthenticated
endpoint.

The name-suffix hash is the detail that makes this better than the
server, not merely equivalent. A ConfigMap named
\code{course-service-config-7b9c4} changes name when its content
changes; the Deployment referencing the new name is a new pod
template; Kubernetes rolls the pods. Config change, rollout, and
rollback are the same operation as a code change --- which is what
Chapter~27's GitOps section wants.

\begin{decision}[when to drop the config-server entirely]
Drop it when three things are true: every runtime is Kubernetes, no
service needs a value to change without a pod restart, and secrets
already live in the platform (they do --- Chapter~16). Keep it when any
is false: services on VMs or in Compose in production, a genuine
runtime-refresh requirement (rare; feature flags are better served by a
purpose-built system), or an organization whose configuration audit
trail must be independent of the deployment repository. An interviewer
who asks ``would you use Spring Cloud Config on Kubernetes?'' is
listening for this list, and for the sentence ``the platform already
provides it''.
\end{decision}

\section{Outside the project}

The git-backed config-server with \code{\{cipher\}} values was the
standard from about 2015 to 2020 and is still what you will find in
most Java shops. Two things have replaced parts of it. HashiCorp Vault
(and its cloud equivalents, AWS Secrets Manager, Azure Key Vault) holds
secrets with leases and rotation; Spring Cloud Vault or the External
Secrets Operator delivers them into pods, and the config repository
holds only references. Feature-flag services (LaunchDarkly, Unleash,
OpenFeature) took over the runtime-refresh use case with targeting
rules and audit that \code{@RefreshScope} never had. What survived
unchanged is the discipline: config outside the image, precedence you
can state from memory, and no working secret in any file you commit.

\section{Questions from real interviews}

\subsection*{Q: You have to rotate a database password across fifty
services with no downtime. Walk me through it.}

Rotation is never a single write; it is an overlap. First, create the
new credential in the database while the old one still works ---
Postgres lets two roles or two passwords coexist, and managed services
(RDS, Cloud SQL) expose ``dual password'' for exactly this. Second, put
the new value into the platform Secret (or the config repository as a
\code{\{cipher\}} value) and trigger the consumers: on Kubernetes that
is a \code{rollout restart} per Deployment, which the ConfigMap
name-hash trick does for you; with Spring Cloud Bus it is one
\code{POST /actuator/busrefresh} that every instance receives over
Kafka. Third, watch the database's connection log until no session uses
the old role, then revoke it. The project's pitfall in Chapter~16 ---
``changed the Secret, nothing happened'' --- is what step two prevents.
The senior answer names the overlap explicitly and says which step is
allowed to fail: revocation, which you simply postpone.

\subsection*{Q: A config change must reach every instance within a
minute. Restarting is too slow. What do you do?}

Distinguish the value kinds first. A log level, a rate limit, a feature
flag: mark the bean \code{@RefreshScope}, expose
\code{/actuator/busrefresh}, add \code{spring-cloud-starter-bus-kafka}
--- the project already has Kafka --- and one POST broadcasts a
\code{RefreshRemoteApplicationEvent}; each instance re-fetches its
document and rebuilds refresh-scoped beans within seconds. A database
URL, a Kafka bootstrap address, a signing key: do \emph{not} refresh
those, because the connection pool, the producer, or the verifier
holds the old value in objects that are not rebuilt cleanly; restart
instead. If the ask is really ``flags for product experiments'', say
that a feature-flag service (Unleash, LaunchDarkly) with an SDK that
polls every few seconds is the purpose-built tool, and
\code{@RefreshScope} is the ten-year-old workaround.

\subsection*{Q: The config server is a single point of failure. How do
you handle that?}

Start by narrowing what it is a SPOF \emph{for}: startup, not traffic.
A running service never calls it again, so a config-server outage
affects only restarts and scale-outs. Then reduce even that: two
replicas behind the Service (it is stateless with a native or git
backend --- a git clone per replica is fine), \code{fail-fast} with
retry so a booting pod waits up to a few minutes instead of dying, and
a readiness probe on the config-server so the Service never routes to
a replica still cloning. The strongest answer removes the dependency:
serve the documents from ConfigMaps, and the ``SPOF'' becomes the
Kubernetes API server, which you already depend on for everything else.

\subsection*{Q: How do you keep secrets out of git while still having
config in git?}

Three layers, each acceptable at a different size. Smallest: the
config-server's \code{/encrypt} with a symmetric key held only by the
server; git stores \code{\{cipher\}} text, and a leaked repository is
useless without the key. Middle: Kubernetes Secrets created out of band
(the project's \code{create-secrets.sh}) and referenced by name in the
manifests, so git holds the shape but never the value. Largest: a
secrets manager --- Vault, AWS Secrets Manager --- with the External
Secrets Operator synchronizing into Kubernetes Secrets, so rotation and
audit live where the secret does. The trap to name: a
\code{\${VAR:default}} with a working default is a secret in git no
matter which layer you chose.

\subsection*{Q: Two pods of the same service are running with different
configuration. How did that happen, and how do you prevent it?}

The usual causes: a rolling update that stalled halfway (old pod on the
old document, new pod on the new one --- normal and temporary); a
config-server whose git clone is stale on one replica because
\code{clone-on-start} was off and one replica never refreshed; or, the
subtle one, a \code{/actuator/refresh} sent to one instance through a
load balancer that happened to pick it. Prevention is to make
configuration part of the pod template: mount it from a ConfigMap
whose name carries a content hash, so ``which config'' is answered by
``which ReplicaSet''. Detection is to expose the config version ---
the git SHA the config-server returns as \code{version}, or the
ConfigMap name --- in \code{/actuator/info}, and to compare across
pods before believing any ``it works on one of them'' report.

\begin{quiz}
\begin{enumerate}
  \item The served \code{course-service.yml} says
  \code{password: root} and the Deployment's Secret says otherwise. State
  which wins and why, then describe what happens if the Secret is
  deleted and the pod restarts.
  \item Explain why \code{\${APP\_S3\_SECRET\_KEY:minioadmin}} in a
  served document is not a precedence rule, and what the config-server
  actually sends to the client for that key.
  \item A colleague runs \code{POST /actuator/refresh} on both
  auth-service and api-gateway after changing \code{JWT\_SECRET}. List
  the two distinct failure windows and name the mechanism from
  Chapter~28 that a hot rotation actually requires.
  \item The config-server dies; course-service is scaled from one to
  two replicas during the outage. Describe what the new pod does with
  \code{configserver:} (non-optional), with \code{optional:configserver:},
  and with \code{fail-fast} plus retry.
  \item After the ConfigMap migration, what causes pods to restart when
  a value changes, and why does that make the migration \emph{more}
  capable than the server it replaced rather than merely equivalent?
  \item Interview: ``Would you deploy Spring Cloud Config on Kubernetes?''
  Give the three conditions under which you would remove it and the
  one condition under which you would keep it.
\end{enumerate}
\end{quiz}
